\documentclass[11pt]{article}
\usepackage{amsmath}
\usepackage{geometry}
\usepackage{setspace}
\usepackage{xcolor}
\usepackage{mdframed}
\usepackage{booktabs}
\usepackage{hyperref}
\usepackage{listings}
\usepackage{array}
\usepackage{longtable}
\geometry{margin=1in}
\setstretch{1.15}

\lstset{
  language=C++,
  basicstyle=\ttfamily\small,
  keywordstyle=\color{blue!70}\bfseries,
  commentstyle=\color{green!50!black}\itshape,
  stringstyle=\color{red!60},
  breaklines=true,
  frame=single,
  backgroundcolor=\color{gray!5},
  numbers=none,
  xleftmargin=6pt,
  xrightmargin=6pt
}

\newcommand{\sth}{\sigma_\theta}
\newcommand{\sal}{\sigma_\alpha}
\newcommand{\sz}{\sigma_z}
\newcommand{\sx}{\sigma_x}
\newcommand{\Lam}{\Lambda}
\newcommand{\ppi}{\pi^-}

\newmdenv[backgroundcolor=blue!5,linecolor=blue!40,linewidth=1pt,
  roundcorner=4pt,innertopmargin=6pt,innerbottommargin=6pt,
  innerleftmargin=8pt,innerrightmargin=8pt]{notebox}
\newmdenv[backgroundcolor=orange!7,linecolor=orange!50,linewidth=1pt,
  roundcorner=4pt,innertopmargin=6pt,innerbottommargin=6pt,
  innerleftmargin=8pt,innerrightmargin=8pt]{warnbox}

\title{\textbf{\texttt{lambda\_analysis.cxx}\\Code Reference Manual}\\[6pt]
       \large EIC $\Lambda \to p + \pi^-$ Analysis --- Far-Forward Region}
\author{C.\ Ayerbe Gayoso, with assistance from Claude (Anthropic)}
\date{\today}

\begin{document}
\maketitle
\tableofcontents
\newpage

%=============================================================
\section{Overview}
%=============================================================

\texttt{lambda\_analysis.cxx} is a ROOT macro that analyses
$\Lambda \to p + \pi^-$ decays in the EIC far-forward region. It reads
Monte Carlo truth from PODIO/EDM4hep files produced by the EIC simulation
framework (dd4hep/GEANT4), fills histograms, and produces three output
PDF files plus one ROOT file containing all histograms. 

Initial version: Dmitry Romanov from the EIC Meson Structure Working Group
\url{ https://github.com/JeffersonLab/meson-structure/blob/main/tutorials/cpp_01_read_edm4eic.cxx}

\subsection{How to run}

Before running this macro, please read the instruction of how to install and run the EIC shell in \url{https://jeffersonlab.github.io/meson-structure/}. 

\textbf{THIS MACRO ONLY RUNS IN THE EIC SHELL}

\begin{lstlisting}
# Interactive ROOT session
root -x -l lambda_analysis.cxx
root [0] lambda_analysis()         // all events, all files
root [0] lambda_analysis(500)      // 500 events per file (quick test)

# Batch mode (no GUI, faster, recommended for full runs)
root -x -l -b -q lambda_analysis.cxx
\end{lstlisting}

\subsection{Dependencies}

\begin{center}
\begin{tabular}{ll}
\toprule
Library & Purpose \\
\midrule
PODIO (\texttt{podioDict}, \texttt{podioRootIO}) & Reading EDM4hep event frames \\
EDM4hep (\texttt{libedm4hepDict}) & MC particle data model \\
\texttt{fmt} & Formatted printing to stdout \\
ROOT (standard) & Histograms, canvases, files, TProfile \\
\bottomrule
\end{tabular}
\end{center}

\subsection{Output files}

\begin{center}
\begin{tabular}{ll}
\toprule
File & Contents \\
\midrule
\texttt{lambda\_analysis\_18x275.root} & All histograms (reopen to replot) \\
\texttt{lambda\_analysis\_18x275.pdf} & Diagnostic + particle analysis plots \\
\texttt{tracker\_study\_18x275.pdf} & Tracker resolution study (12 pages) \\
\bottomrule
\end{tabular}
\end{center}

%=============================================================
\section{Configuration Namespaces}
%=============================================================

All tunable parameters live in two namespaces at the top of the file.
\textbf{No other part of the code needs to be changed to adjust cuts or
tracker geometry.}

\subsection{\texttt{namespace Cuts}}

\begin{center}
\renewcommand{\arraystretch}{1.3}
\begin{tabular}{lll}
\toprule
Constant & Value & Meaning \\
\midrule
\texttt{LAMBDA\_PDG} & 3122 & PDG code of the $\Lambda$ baryon \\
\texttt{LAMBDA\_PRIMARY\_IDX} & 4 & MC particle index of primary $\Lambda$ \\
\texttt{LAMBDA\_N\_DAUGHTERS} & 2 & Selects charged decay ($p + \pi^-$) \\
\texttt{DECAY\_Z\_MIN} & 20800\,mm & Last beamline magnet exit \\
\texttt{DECAY\_Z\_MAX} & 35000\,mm & ZDC face \\
\texttt{ZDC\_Z} & 35000\,mm & ZDC face (used for projection) \\
\texttt{CROSSING\_ANGLE} & 25\,mrad & Hadron beam crossing angle \\
\texttt{PDG\_PROTON} & 2212 & PDG code of proton \\
\texttt{PDG\_PION\_MINUS} & $-211$ & PDG code of $\pi^-$ \\
\texttt{PDG\_NEUTRON} & 2112 & PDG code of neutron \\
\texttt{PDG\_PION\_ZERO} & 111 & PDG code of $\pi^0$ \\
\texttt{N\_FILES} & 20 & Number of input ROOT files \\
\texttt{FILE\_PATTERN} & path template & \texttt{fmt::format} pattern \\
\texttt{OUTPUT\_PDF} & string & Base name for analysis PDF \\
\bottomrule
\end{tabular}
\end{center}

\begin{warnbox}
\texttt{LAMBDA\_PRIMARY\_IDX = 4} is generator-dependent. If the
simulation configuration changes, this index must be updated. Verify
by running with a small sample and checking the \texttt{[NON-PRIMARY]}
diagnostic output (see Section~\ref{sec:counters}).
\end{warnbox}

\subsection{\texttt{namespace Tracker}}

\begin{center}
\renewcommand{\arraystretch}{1.3}
\begin{tabular}{lll}
\toprule
Constant & Value & Meaning \\
\midrule
\texttt{N\_PLANES} & 4 & Number of tracking planes \\
\texttt{LENGTH} & 2000\,mm & Total tracker length (plane 1 to plane 4) \\
\texttt{STANDOFF} & 500\,mm & Gap from ZDC face to last plane \\
\texttt{SIGMA\_X[]} & 0.05--1.0\,mm & Five pitch scenarios \\
\texttt{N\_SC} & 5 & Number of scenarios \\
\texttt{LABELS[]} & strings & ROOT title labels per scenario \\
\texttt{EFF\_THRESH[]} & 500, 1000, 2000\,mm & Efficiency thresholds \\
\texttt{N\_THRESH} & 3 & Number of thresholds \\
\texttt{PRIMARY\_THRESH} & 1 & Index of primary threshold (1000\,mm) \\
\texttt{OUTPUT\_PDF} & string & Base name for tracker study PDF \\
\bottomrule
\end{tabular}
\end{center}

%=============================================================
\section{Global State}
%=============================================================

\begin{description}
  \item[\texttt{map<string, TH1*> H}] --- Global histogram registry.
        All histograms are stored here by name. Access via the typed
        helpers \texttt{h1()} and \texttt{h2()} (see Section~\ref{sec:histmap}).

  \item[\texttt{string current\_file}] --- Set to the filename currently
        being processed by \texttt{process\_file()}. Used in diagnostic
        print statements to identify which file contains anomalous events.

  \item[\texttt{Counters counters}] --- Global instance of the
        \texttt{Counters} struct. Incremented throughout
        \texttt{process\_event()} and printed at the end of the run.
\end{description}

%=============================================================
\section{Global Functions (Physics Utilities)}
%=============================================================

\subsection{\texttt{sigma\_theta\_rad(sigma\_x\_mm)}}

Returns angular resolution [rad] assuming all $N=4$ planes resolve the
tracks (best case). Retained for documentation and the summary table.
In the event loop, use \texttt{sigma\_theta\_eff()} instead.

\subsection{\texttt{plane\_z(ip)}}

Returns the absolute $z$ position [mm] of tracker plane \texttt{ip}
(1-indexed). Planes are equally spaced:
\[
  z_1 = 32{,}500\,\text{mm}, \quad z_4 = 34{,}500\,\text{mm}
\]

\subsection{\texttt{sigma\_theta\_eff(sigma\_x\_mm, first\_plane)}}

Returns the effective angular resolution [rad] using only planes from
\texttt{first\_plane} to \texttt{N\_PLANES}. Returns $10^9$ (sentinel
= unresolvable) when fewer than 2 planes are available.

\begin{center}
\begin{tabular}{cccc}
\toprule
\texttt{first\_plane} & $N_\text{eff}$ & $L_\text{eff}$ (mm)
    & Relative $\sth$ \\
\midrule
1 & 4 & 2000 & 1.0 (best) \\
2 & 3 & 1333 & 1.15 \\
3 & 2 & 667  & 1.58 \\
4 & 1 & 0    & $\infty$ \\
\bottomrule
\end{tabular}
\end{center}

\subsection{\texttt{PolarAngle(px, py, pz)}}
Returns $\theta = \arccos(p_z/|\vec{p}|)$ in degrees.

\subsection{\texttt{PhiAngle(px, py)}}
Returns $\phi = \text{atan2}(p_y, p_x)$ in degrees.

\subsection{\texttt{OpeningAngle(a, b)}}
Returns the opening angle between two \texttt{DaughterKin} objects in
degrees, from the dot product of unit momentum vectors.

%=============================================================
\section{Data Structures}
%=============================================================

\subsection{\texttt{struct DaughterKin}}

Holds kinematic quantities for one decay daughter projected onto the ZDC
plane. Computed by \texttt{compute()}.

\begin{center}
\renewcommand{\arraystretch}{1.3}
\begin{tabular}{lll}
\toprule
Member & Unit & Meaning \\
\midrule
\texttt{theta} & degrees & Polar angle w.r.t.\ $z$-axis \\
\texttt{phi} & degrees & Azimuthal angle \\
\texttt{xAtZDC}, \texttt{yAtZDC} & mm & Position at ZDC face (lab frame) \\
\texttt{radius} & mm & $\sqrt{x^2+y^2}$ in lab frame \\
\texttt{xBC}, \texttt{yBC} & mm & Position relative to hadron beam axis \\
\texttt{radiusBC} & mm & Distance from beam axis \\
\bottomrule
\end{tabular}
\end{center}

The vertex offset (crossing angle) is applied inside \texttt{compute()}
using the true decay vertex position $(x_v, y_v, z_v)$.
\texttt{isSet()} returns \texttt{true} if \texttt{theta != 0}.

\subsection{\texttt{struct Counters}}
\label{sec:counters}

Tracks particle counts at each selection stage. Printed to stdout by
\texttt{counters.print()} at the end of the run.

\begin{center}
\renewcommand{\arraystretch}{1.3}
\begin{tabular}{ll}
\toprule
Member & Counts \\
\midrule
\texttt{total\_events} & Events read from all files \\
\texttt{lambda\_escaped} & Primary $\Lambda$ with no daughters \\
\texttt{lambda\_before\_magnet} & Primary $\Lambda$ decaying before 20.8\,m \\
\texttt{lambda\_beyond\_zdc} & Primary $\Lambda$ decaying past 35\,m \\
\texttt{lambda\_field\_free} & Primary $\Lambda$ in field-free window (any channel) \\
\texttt{lambda\_charged} & Field-free, charged channel $p+\pi^-$ [used] \\
\texttt{lambda\_neutral} & Field-free, neutral channel $n+\pi^0$ [not used] \\
\texttt{lambda\_other\_ndau} & Field-free, material interactions ($N_\text{dau}>2$) \\
\texttt{proton} & Protons found in charged channel \\
\texttt{pion\_minus} & $\pi^-$ found in charged channel \\
\texttt{neutron} & Neutrons (neutral channel, informational) \\
\texttt{pion\_zero} & $\pi^0$ (neutral channel, informational) \\
\bottomrule
\end{tabular}
\end{center}

%=============================================================
\section{Histogram Registry}
\label{sec:histmap}
%=============================================================

\subsection{Design}

All histograms live in \texttt{map<string, TH1*> H}. Two typed accessors
avoid casting at every fill site:
\begin{lstlisting}
inline TH1F* h1(const string& name) { return static_cast<TH1F*>(H.at(name)); }
inline TH2F* h2(const string& name) { return static_cast<TH2F*>(H.at(name)); }
\end{lstlisting}

For \texttt{TProfile} histograms, cast explicitly at the fill site:
\begin{lstlisting}
static_cast<TProfile*>(H.at("ResZ_rms_vsAlpha_sc0"))->Fill(alpha, dz);
\end{lstlisting}

\subsection{How to add a new histogram}
\label{sec:add_histogram}

\begin{enumerate}
  \item In \texttt{book\_histograms()}: add one line:
\begin{lstlisting}
H["MyHist"] = new TH1F("MyHist", "Title;X;Y", 100, 0, 100);
\end{lstlisting}
  \item In \texttt{process\_event()}: fill it where the quantity is computed:
\begin{lstlisting}
h1("MyHist")->Fill(my_value);
\end{lstlisting}
  \item In \texttt{plot()} or \texttt{plot\_tracker\_study()}: draw it.
  \item It is automatically written to the ROOT file by
        \texttt{save\_histograms()} --- no extra step needed.
\end{enumerate}

\subsection{Complete histogram table}

\subsubsection*{Lambda kinematics}
\begin{longtable}{p{4.5cm}lp{6cm}}
\toprule
Key & Type & Contents \\
\midrule
\texttt{LambdaEnd} & TH1F & Decay endpoint $z$, all $\Lambda$ with 2 parents [0--35\,m] \\
\texttt{LambdaEnd\_1} & TH1F & Decay endpoint [0--6\,m] \\
\texttt{LambdaEnd\_2} & TH1F & Decay endpoint [6--20.8\,m] \\
\texttt{LambdaEnd\_3} & TH1F & Decay endpoint, primary $\Lambda$, field-free, any channel \\
\texttt{LambdaEnd\_analysis} & TH1F & Decay endpoint, primary, field-free, charged channel only.
                                      Integral = \texttt{counters.lambda\_charged} \\
\texttt{LambdaGenZ} & TH1F & Production vertex $z$ \\
\texttt{LambdaZnoDaughters} & TH1F & Endpoint of escaped $\Lambda$ (no daughters) \\
\texttt{LambdaAngleGen} & TH1F & Polar angle at production \\
\bottomrule
\end{longtable}

\subsubsection*{Proton daughters}
\texttt{Vertex\_p}, \texttt{Endpoint\_p}, \texttt{Theta\_p},
\texttt{XYatZDC\_p} (TH2F, lab), \texttt{XYatZDC\_p\_BC} (TH2F, beam-centred),
\texttt{ZvsTheta\_p}, \texttt{ZvsRadius\_p}.

\subsubsection*{Pion daughters}
Same set with \texttt{\_pi} suffix.

\subsubsection*{Neutral channel (informational)}
\texttt{Vertex\_n}, \texttt{Endpoint\_n}, \texttt{Theta\_n},
\texttt{XYatZDC\_n} for neutrons; same with \texttt{\_pi0} for $\pi^0$.

\subsubsection*{Correlations}
\texttt{ThetaCorr} (TH2F), \texttt{RadiusSumVsZ} (TH2F),
\texttt{DeltaR}, \texttt{OpenAngle}, \texttt{DeltaRvsZ} (TH2F).

\subsubsection*{Tracker study: $\Delta R$ at planes (truth)}
\begin{longtable}{p{4.5cm}lp{6cm}}
\toprule
Key & Type & Contents \\
\midrule
\texttt{DeltaR\_pl\{1-4\}} & TH1F & Track separation at plane $i$, full range [0--250\,mm] \\
\texttt{DeltaR\_pl\{1-4\}\_zoom} & TH1F & Track separation at plane $i$, zoomed [0--3\,mm],
                                           0.05\,mm/bin. Shows sub-mm region where pitch matters \\
\bottomrule
\end{longtable}

\subsubsection*{Tracker study: analytical (suffix \texttt{\_sc0} to \texttt{\_sc4})}
\begin{longtable}{p{4.5cm}lp{6cm}}
\toprule
Key & Type & Contents \\
\midrule
\texttt{SigmaZ\_sc\{i\}} & TH1F & $\sigma_z$ distribution [0--5000\,mm] \\
\texttt{SigmaAlpha\_sc\{i\}} & TH1F & $\sigma_\alpha$ distribution [0--5\,mrad] \\
\texttt{SigmaZ\_vsAlpha\_sc\{i\}} & TH2F & $\sigma_z$ vs true $\alpha$ [0--15\,mrad] \\
\texttt{SigmaZ\_vsD\_sc\{i\}} & TH2F & $\sigma_z$ vs $D_\text{eff}$ \\
\bottomrule
\end{longtable}

\subsubsection*{Tracker study: smearing (suffix \texttt{\_sc0} to \texttt{\_sc4})}
\begin{longtable}{p{4.8cm}lp{5.7cm}}
\toprule
Key & Type & Contents \\
\midrule
\texttt{ResZ\_sc\{i\}} & TH1F & $z_\text{reco} - z_\text{true}$ residual \\
\texttt{ResAlpha\_sc\{i\}} & TH1F & $\alpha_\text{reco} - \alpha_\text{true}$ residual \\
\texttt{ResZ\_vsAlpha\_sc\{i\}} & TH2F & $\Delta z$ vs true $\alpha$ [0--15\,mrad].
                                         Use \texttt{ProjectionY(name,b,b)} to get
                                         the $\Delta z$ distribution per $\alpha$ bin \\
\texttt{ZReco\_sc\{i\}} & TH1F & Reconstructed vertex $z$ [15--40\,m] \\
\texttt{SmearEff\_thr\{j\}\_sc\{i\}} & TH1F & Pass/fail histogram for
                                              $|\Delta z| < \texttt{EFF\_THRESH[j]}$.
                                              \texttt{GetBinContent(2)/GetEntries()}
                                              = precision efficiency \\
\texttt{ResZ\_rms\_vsAlpha\_sc\{i\}} & TProfile & TProfile(``S'') of $\Delta z$ vs $\alpha$.
                                                   Filled but not used for plotting;
                                                   plotting uses \texttt{ProjectionY}
                                                   of \texttt{ResZ\_vsAlpha} instead \\
\bottomrule
\end{longtable}

%=============================================================
\section{Functions Reference}
%=============================================================

\subsection{\texttt{book\_histograms()}}

Called once at startup. Creates every histogram in \texttt{H}. All
booking is here --- change ranges or binning only in this function.

\subsection{\texttt{process\_event(event, event\_number)}}

Main event processing function. Called once per MC event.

\textbf{Execution flow:}
\begin{enumerate}
  \item Retrieve \texttt{MCParticles} from the PODIO frame.
  \item Loop over all particles. For each $\Lambda$ with 2 parents:
        fill \texttt{LambdaEnd} histograms and \texttt{LambdaAngleGen}.
        If the particle is in the field-free window and is not the
        primary Lambda (index $\neq 4$), print a diagnostic.
  \item For all particles: fill \texttt{LambdaGenZ}.
  \item Select primary Lambda (index = \texttt{LAMBDA\_PRIMARY\_IDX}).
        If no daughters: fill \texttt{LambdaZnoDaughters},
        increment \texttt{lambda\_escaped}, skip.
  \item Apply $z$ window cut. Increment
        \texttt{lambda\_before\_magnet} or \texttt{lambda\_beyond\_zdc}
        accordingly.
  \item Increment \texttt{lambda\_field\_free}.
  \item Check daughter PDGs to identify channel. Increment
        \texttt{lambda\_neutral} + fill neutral histograms + skip, or
        \texttt{lambda\_other\_ndau} + skip, as appropriate.
  \item For the charged channel: increment \texttt{lambda\_charged},
        fill \texttt{LambdaEnd\_analysis}.
  \item Daughter loop: compute \texttt{DaughterKin} for $p$ and $\pi^-$,
        fill all daughter histograms.
  \item If both daughters found: fill correlation histograms
        (\texttt{DeltaR}, \texttt{OpenAngle}, etc.).
  \item Tracker study (if $\alpha > 10^{-6}$\,rad):
        \begin{enumerate}
          \item Fill \texttt{DeltaR\_pl\{i\}} and \texttt{DeltaR\_pl\{i\}\_zoom}
                for each plane.
          \item For each scenario: find \texttt{first\_plane} where
                $\Delta R > \sx$. If none: skip.
          \item Compute \texttt{sigma\_theta\_eff}, $D_\text{eff}$, $\sigma_z$.
                Fill analytical histograms.
          \item Smear slopes with $\mathcal{G}(0, \sigma_\theta)$.
                Reconstruct $z_\text{reco}$ via PCA. Set \texttt{reco\_ok}
                flag only when the PCA succeeded (not the fallback).
          \item Fill \texttt{ResZ}, \texttt{ResAlpha}, \texttt{ResZ\_vsAlpha},
                \texttt{ZReco}. If \texttt{reco\_ok}: also fill
                \texttt{SmearEff\_thr\{j\}} and \texttt{ResZ\_rms\_vsAlpha}.
        \end{enumerate}
\end{enumerate}

\subsection{\texttt{process\_file(filename, events\_limit)}}

Sets \texttt{current\_file = filename}, opens one PODIO ROOT file,
calls \texttt{process\_event()} for each entry up to \texttt{events\_limit}
($-1$ = all). Wrapped in try/catch.

\subsection{\texttt{plot()}}

Creates diagnostic and particle analysis canvases. Returns a
\texttt{vector<TCanvas*>} for \texttt{save\_plots()}.

\subsection{\texttt{save\_plots(pages, output\_name)}}

Writes a canvas vector to a multi-page PDF using ROOT's
\texttt{Print("name.pdf(")} / \texttt{Print("name.pdf")} /
\texttt{Print("name.pdf)")} pattern.

\subsection{\texttt{plot\_tracker\_study()}}

Creates all 12 tracker study canvases. Returns a
\texttt{vector<TCanvas*>} for \texttt{save\_tracker\_study()}.

The 12 pages in order:
\begin{enumerate}
  \item $\Delta R$ at each tracker plane (2$\times$2)
  \item $\Delta R$ zoomed overlay [0--3\,mm], all planes
  \item $\sigma_z$ distributions overlay (analytical)
  \item $\sigma_\alpha$ distributions overlay (analytical)
  \item $\sigma_z$ vs $\alpha$ per scenario (3$\times$2)
  \item $\sigma_z$ vs $D_\text{eff}$ per scenario (3$\times$2)
  \item Summary: mean $\sigma_z$ + analytical efficiency curves
  \item $z_\text{reco}$ distributions + precision smearing efficiency
  \item $\Delta z$ residuals overlay
  \item $\Delta\alpha$ residuals overlay
  \item $\Delta z$ vs $\alpha$ per scenario (3$\times$2)
  \item Validation: RMS($\Delta z$) vs analytical $\sigma_z$ per $\alpha$ bin
\end{enumerate}

\subsection{\texttt{save\_tracker\_study(pages)}}

Writes the tracker canvas vector to \texttt{tracker\_study\_18x275.pdf}.

\subsection{\texttt{save\_histograms(rootfile)}}

Opens a ROOT file in RECREATE mode and writes every histogram in
\texttt{H} with \texttt{->Write()}. Called automatically from
\texttt{lambda\_analysis()}. To reopen later:
\begin{lstlisting}
TFile* f = TFile::Open("lambda_analysis_18x275.root");
f->ls();                          // list all histograms
TH1F* h = (TH1F*)f->Get("DeltaR");
h->Draw();
\end{lstlisting}

\subsection{\texttt{lambda\_analysis(events\_limit = -1)}}

Entry point. Calls in order:
\begin{enumerate}
  \item \texttt{book\_histograms()}
  \item \texttt{process\_file()} for each of the \texttt{N\_FILES} input files
  \item \texttt{counters.print()}
  \item \texttt{save\_histograms("lambda\_analysis\_18x275.root")}
  \item \texttt{plot()} + \texttt{save\_plots()} $\to$ analysis PDF
  \item \texttt{plot\_tracker\_study()} + \texttt{save\_tracker\_study()} $\to$ tracker PDF
\end{enumerate}

%=============================================================
\section{Event Selection Logic}
%=============================================================

Applied in \texttt{process\_event()} in order:

\begin{center}
\renewcommand{\arraystretch}{1.4}
\begin{tabular}{lp{9cm}}
\toprule
Cut & Reason \\
\midrule
$\texttt{PDG} = 3122$ & Select $\Lambda$ baryons \\
2 parents & Select from primary interaction vertex \\
Index $= 4$ & Select primary $\Lambda$ from hard scatter \\
$\texttt{getDaughters().size()} > 0$ & Lambda must have decayed \\
$z_v \in (20800, 35000)\,\text{mm}$ & Field-free decay window \\
Exactly 2 daughters & Charged channel only \\
Daughters are $p + \pi^-$ & Confirmed by PDG code check \\
\bottomrule
\end{tabular}
\end{center}

\subsection{How to add a new cut}

\begin{enumerate}
  \item Add a counter to \texttt{struct Counters} for the rejected events.
  \item In \texttt{process\_event()}, add the cut condition with
        \texttt{continue} to skip the event.
  \item Increment the counter before \texttt{continue}.
  \item Add the counter to \texttt{Counters::print()}.
\end{enumerate}

Example --- adding a minimum $\alpha$ cut of 2\,mrad:
\begin{lstlisting}
// In struct Counters:
int lambda_small_alpha = 0;

// In process_event(), after computing alpha:
if (alpha * M_PI / 180.0 < 2e-3) {
    counters.lambda_small_alpha++;
    continue;
}

// In Counters::print():
fmt::print("    Small alpha (< 2 mrad)   : {:6d}\n", lambda_small_alpha);
\end{lstlisting}

%=============================================================
\section{How to Modify the Analysis}
%=============================================================

\subsection{Add a new tracker pitch scenario}

\begin{enumerate}
  \item Add the pitch value to \texttt{Tracker::SIGMA\_X[]}
  \item Increment \texttt{Tracker::N\_SC} by 1
  \item Add a label to \texttt{Tracker::LABELS[]}
\end{enumerate}
All histogram booking, filling, and plotting loops use \texttt{N\_SC}
and update automatically.

\subsection{Change the primary efficiency threshold}

Change \texttt{Tracker::PRIMARY\_THRESH} to the index of the desired
threshold in \texttt{EFF\_THRESH[]}. The summary plot and comparison
table update automatically.

\subsection{Change the decay region}

Edit \texttt{Cuts::DECAY\_Z\_MIN} and \texttt{Cuts::DECAY\_Z\_MAX}.
All cuts, histograms, and efficiency calculations use these constants.

\subsection{Process a different beam energy}

Change \texttt{Cuts::N\_FILES}, \texttt{Cuts::FILE\_PATTERN}, and
\texttt{Tracker::OUTPUT\_PDF}. Run \texttt{lambda\_analysis()} and
the new files will be processed.

\subsection{Quick test run}
\begin{lstlisting}
root [0] lambda_analysis(100)   // 100 events per file
\end{lstlisting}

%=============================================================
\section{How Histograms Are Saved}
%=============================================================

\subsection{To ROOT file}

\texttt{save\_histograms()} iterates over the entire map \texttt{H}
and calls \texttt{->Write()} on every entry. This happens automatically
--- no histogram needs to be explicitly listed. Any histogram booked in
\texttt{book\_histograms()} is saved.

\subsection{To PDF}

Both \texttt{save\_plots()} and \texttt{save\_tracker\_study()} use the
same pattern:
\begin{lstlisting}
// Open the PDF
pages[0]->Print("output.pdf(");
// Append middle pages
for (int i = 1; i < (int)pages.size()-1; ++i)
    pages[i]->Print("output.pdf");
// Close the PDF
pages.back()->Print("output.pdf)");
\end{lstlisting}
Canvases are accumulated in a \texttt{vector<TCanvas*>} inside
\texttt{plot()} or \texttt{plot\_tracker\_study()} and passed to the
save function.

To add a new page to the tracker PDF, add a new canvas block in
\texttt{plot\_tracker\_study()} before \texttt{return pages;} and call
\texttt{pages.push\_back(c)} at the end of the block.

%=============================================================
\section{Known Limitations}
%=============================================================

\begin{itemize}
  \item \textbf{No background.} All efficiencies are upper bounds.
  \item \textbf{No momentum measurement.} The tracker alone cannot
        reconstruct the $\Lambda$ invariant mass.
  \item \textbf{No multiple scattering.} Thin silicon assumed.
  \item \textbf{Conservative $N_\text{eff}$ model.} Events where tracks
        first separate at the last plane are treated as unresolvable.
  \item \textbf{\texttt{LAMBDA\_PRIMARY\_IDX = 4} is generator-dependent.}
        Verify for new simulation configurations.
  \item \textbf{Smearing uses Gaussian noise only.} Real detector effects
        (dead channels, misalignment, hadronic interactions in material)
        are not modelled.
\end{itemize}

\end{document}
